\documentclass[10pt,a4paper]{amsart}
\usepackage[margin=19mm]{geometry}
\usepackage{amsmath,amssymb,mathtools}
\usepackage[scr=rsfs,cal=euler]{mathalfa}
\usepackage{xfrac}
\usepackage[unicode,hidelinks,bookmarks=false]{hyperref}
\newcommand{\ie}{i.e.\ }
\newcommand{\cf}{cf.\ }
\newcommand{\resp}{respectively\ }
\newcommand{\Subsection}{\subsection}
\providecommand{\nobreakdash}{\nobreak\hbox{-}}
\setlength{\emergencystretch}{2em}
\multlinegap0pt
\usepackage{amssymb}
\usepackage{enumerate}
\usepackage{tikz}
\usepackage{braket}
\usepackage{caption}
\newtheorem{corollary}{Corollary}
\newtheorem{lemma}{Lemma}
\renewcommand{\H}{\mathcal{H}}
\newcommand{\RR}{\mathbb{R}}
\newcommand{\SU}{\operatorname{SU}}
\newcommand{\Sp}{\operatorname{Sp}}
\newcommand{\U}{\operatorname{U}}
\renewcommand{\a}{\mathfrak{a}}
\newcommand{\g}{\mathfrak{g}}
\newcommand{\h}{\mathfrak{h}}
\renewcommand{\k}{\mathfrak{k}}
\newcommand{\p}{\mathfrak{p}}
\renewcommand{\sp}{\mathfrak{sp}}
\newcommand{\su}{\mathfrak{su}}
\title{Real 3-qubit gate decompositions via triality}
\author{Brendan Pawlowski}
\date{Source: arXiv:2512.18049v2 (CC BY 4.0)}
\begin{document}
\maketitle

\noindent\emph{Source and licence.} Real 3-qubit gate decompositions via triality. Version arXiv:2512.18049v2 (CC BY 4.0), distributed by arXiv under the Creative Commons Attribution 4.0 International licence, http://creativecommons.org/licenses/by/4.0/, as recorded in the arXiv licence metadata for this exact version and shown on the abstract page https://arxiv.org/abs/2512.18049v2. Full licence notice: \texttt{licenses/arxiv-2512.18049-CC-BY-4.0.txt}.\par\smallskip

\noindent\textit{Editorial scope.} Counted unit: the article's lemma lem:real-gates-sp (source0.tex lines 1076--1088). The article's complete original proof of it is delivered with it (source0.tex lines 1090--1146, one \texttt{proof} environment of 57 lines). No mathematics is written, repaired or re-derived: the statement and the proof are the article's own lines, in the article's own notation, and no retained line is substituted or re-flowed.  Retained uncounted as the source-local context the proof needs: lines 225--227.  Nothing was omitted from any retained span, so the package carries no omission record.  No retained line contains a citation, so the wrapper carries no bibliography and no bibliography entry.  No retained line contains a figure environment, an image-inclusion command or drawing code, so no image is delivered and the package declares no asset dependency.  Editorial formatting only, in addition: an AMS \texttt{amsart} preamble that restates the article's own declarations and macros used by the retained lines, namely the environments \texttt{corollary}, \texttt{lemma} on separate counters, together with the standard packages those lines need.  The retained environments print in this document as Corollary 1, Lemma 1 in order of appearance, whereas the article prints the same items with the numbering of its own full document.  The complete native source of the article as distributed in the arXiv e-print for this exact version is bundled unchanged as \texttt{sources/191359/source0.tex}.
\begin{corollary} \label{cor:ka}
$\g$ is generated as a Lie algebra by $\k$ and $\a$.
\end{corollary}

\begin{lemma} \label{lem:real-gates-sp} The subgroup of $\U(4)$ generated by
\begin{equation*}
\H_1 = R_y(\RR) \otimes \SU(2) \quad \text{and} \quad \H_2 = C_1^2 (R_x(\RR) \otimes \SU(2))C_1^2
\end{equation*}
is
\begin{equation*}
\Sp(2) = \left\{\left[\begin{smallmatrix} a & -\overline{b} & c & -\overline{d} \\
								   b & \overline{a} & d & \overline{c}\\
								   e & -\overline{f} & g & -\overline{h}\\
							       f & \overline{e} & h & \overline{g}
\end{smallmatrix}\right] \in \U(4) \right\}
\end{equation*}
\end{lemma}

\begin{proof}
It suffices to show that the Lie algebras
\begin{align*}
\h_1 &= i \RR\sigma_y \otimes I_2 + I_2 \otimes \su(2) = \left\{ \left[ \begin{smallmatrix}
ir & -\overline{a} & -s & 0\\
a &  -ir & 0 & -s\\
s & 0 & ir & -\overline{a}\\
0 & s & a & -ir \end{smallmatrix} \right] : r,s \in \RR\right\} \\
\h_2 &= C_1^2 (i \RR \sigma_x \otimes I_2 + I_2 \otimes \su(2)) C_1^2 = \left\{ \left[ \begin{smallmatrix}
ir & -\overline{a} & 0 & is\\
a &  -ir & is & 0\\
0 & is & -ir & a\\
is & 0 & -\overline{a} & ir
\end{smallmatrix} \right] : r,s \in \RR\right\}
\end{align*}
generate the Lie algebra
\begin{equation*}
\sp(2) = \left\{ \left[ \begin{smallmatrix}
ir & -\overline{a} & -b & -\overline{c}\\
a &  -ir 		  & c   & -\overline{b}\\
b & -\overline{c}  & is & -\overline{d}\\
c & \overline{b}   & d  & -is
\end{smallmatrix}\right] : r,s \in \RR \right\}.
\end{equation*}

Let $\theta : \Sp(2) \to \Sp(2)$ be conjugation by $\sigma_x \otimes \sigma_x$, so $\theta$ (and $d\theta$) have the effect of rotating a matrix by $180^\circ$. The subspace $\k \subseteq \sp(2)$ fixed by $d\theta$ is exactly $\h_2$. 

The $(-1)$-eigenspace of $d\theta$ is
\begin{equation} \label{eq:sp2-cartan}
\p = \left\{ \left[ \begin{smallmatrix}
ir & -\overline{a} & -\overline{b} & -s\\
a &  -ir 		  & s   & -b\\
b & -s  & ir & -a\\
s & \overline{b}   & \overline{a}  & -ir
\end{smallmatrix}\right] : r,s \in \RR \right\}.
\end{equation}
Let $\a$ be the span of the two commuting elements
\begin{equation*}
i\sigma_y \otimes I_2 = \left[ \begin{smallmatrix}
0 & 0 & 1 & 0\\
0 & 0 & 0 & 1\\
-1 &0 & 0 & 0\\
0 & -1 & 0 & 0
\end{smallmatrix} \right] \qquad \text{and} \qquad I_2 \otimes i\sigma_y = \left[ \begin{smallmatrix}
0 & 1 & 0 & 0\\
-1 & 0 & 0 & 0\\
0 & 0 & 0 & 1\\
0 & 0 & -1 & 0
\end{smallmatrix} \right].
\end{equation*}
Note that $\a \subseteq \h_1$ and that $\a$ is also contained in the $(-1)$-eigenspace $\p$ of $d\theta$. We claim $\a$ is a maximal abelian subalgebra of $\p$. Indeed, a matrix $A$ commutes with $I_2 \otimes i\sigma_y$ if and only if each $2 \times 2$ block has the form $\left[ \begin{smallmatrix} x & -y \\ y & x \end{smallmatrix} \right]$; in particular, if $A$ has the form \eqref{eq:sp2-cartan} then we must have $r = 0$ and $a,b \in \RR$. Similarly, $A$ commutes with $i\sigma_y \otimes I_2$ iff $a,b \in \RR$ and $s = 0$. Thus if both conditions hold then $A =  \left[ \begin{smallmatrix}
0 & -a & -b & 0\\
a &  0 		  & 0   & -b\\
b & 0  & 0 & -a\\
0 & b   & a  & 0
\end{smallmatrix}\right]$, which already lies in $\a$. Cartan decomposition (specifically Corollary~\ref{cor:ka}) therefore shows $\sp(2)$ is generated by $\k = \h_2$ together with $\a \subseteq \h_1$.
\end{proof}

\end{document}
